\section{Methodik}
\label{sec:3}

Der folgende Abschnitt gliedert sich in drei Teile. Zunächst wird die Datengrundlage dieser Arbeit - das \textit{konkorp} - vorgestellt. Anschließend erfolgt die Schilderung der quantitativen und zuletzt der qualitativen Untersuchung.

\subsection{Datengrundlage}
 konkorp dies das

\subsection{Annotationen}

Zunächst wird für ersteren das \textit{konkorp} in ANNIS eingelesen und mittels regulärer Ausdrücke nach (nach syntaktischer Definition) Verben durchsucht. Es wird eine Frequenzanalyse von Verbphrasen durchgeführt. Dies geschieht für jeden der sechs Datensätze separat, um die Ergebnisse auch diachron auswerten zu können.

Es werden schließlich qualitative Analysen einzelner Verbokkurenzen vorgenommen. Hierzu wird auf die Dreiteilung des Transitivitätskonzeptes (processes, participants, circumstances) nach Halliday zurückgegriffen.
Halliday schreibt Sätzen (\textit{clauses}) drei Metaebenen von Bedeutung zu; Austausch, Nachricht, Repräsentation. Die Analyse wird sich vorranging auf die letzte Ebene beziehen, denn auf dieser findet die Konzeption von Transitivität in Hallidays Framework statt.

Auf strukturalistischer Ebene korrelieren Komponenten von Transititivät mit syntaktischen Einheiten (process -- Verbphrase, participants -- Nominalphrase, circumstances -- Präpositionalphrase und Adverbien). Da die Annotation von Konstituenten im \textit{konkorp} bereits erfolgt ist, muss hier keine Analyse mehr erfolgen.
Vielmehr wird das Korpus um zwei Annotationsebenen erweitert; Verbkomplexität und \textit{process type}. In ersterer wird annotiert, ob es sich um ein Verbum simplex, ein Präfix- oder ein Partikelverb handelt. In zweiterer werden die Verben einer der sechs Kategorien nach \citet[171]{Halliday2004} zugeordnet und jeweils die Umstände und Beteiligten des durch das Verb beschriebenen Vorganges identifiziert. Dies soll für alle sechs zugrundeliegenden Verfassungen geschehen. So sollen die einzelnen Verfassungstexte anhand der Nutzung der jeweiligen Verben charakterisiert werden. Ziel ist es, die Hypothese zu überprüfen, ob die Verfassungen mit zunehmender Aktualität einen Wandel in der Dynamik zwischen Staat und Bürger*innen verzeichnen lassen. Hierzu wird in der Analyse und Diskussion betrachtet, welche \textit{process types} in welchem Datensatz wie häufig vorhanden sind und ob sich durch die jeweils in der \textit{clause} enthaltenen \textit{participants} und \textit{circumstances} Rückschlüsse auf verfassungsrechtliche Konzepte ziehen lassen.

\subsection{Qualitative Analyse}

So könnte beispielsweise untersucht werden, welche Rolle den Staatsbürger*innen und welche dem Staat zugeordnet werden können. Welches Organ hat in welcher Epoche welche Rechte und Pflichten, wer handelt, und wer wird behandelt?

Im Idealfall kann linguistische Evidenz dafür gefunden werden, dass aktuellere Verfassungstexte dem Schutz und dem Einräumen von Rechten gegenüber den Bürger*innen dienen, während ältere Verfassungen das Ziel hatten, die Staatsgewalt und ihre Macht zu sichern.

